%% Your notation and commands: derivatives, vectors, operators, sets, dimensionless numbers. Change them to
%% your field's conventions, and add your own here. (The template's own commands, such as \repeatcaption,
%% \doi and \chapterbasedon, are in template/commands.tex.)

%% Derivatives and differentials with the derivative package (loaded by the template), upright d and D:
%%   \odv{y}{x}            dy/dx              \odv[order=2]{y}{x}     d^2y/dx^2 (also \odv[2]{y}{x})
%%   \pdv{u}{t}            partial            \pdv[order={2,1}]{u}{x,y}  mixed: d^3u/dx^2 dy
%%   \mdv{\vec{u}}{t}      material: Du/Dt    \odv*{}{x} f            the function after the fraction
%%   \odif{x}  \pdif{x}  \mdif{x}             the differentials dx, partial x, Dx, as in an integral
%% Compact derivatives, the template's own: \cpdv{u}{t} gives \partial_t u, \cpdv[2]{u}{x} \partial^2_x u.
%% The scripts attach straight to the d or \partial, so the subscript sits close to it.
\newcommand\ThesisCompactDeriv[4]{#1\if\relax\detokenize{#2}\relax\else^{#2}\fi_{#4}{#3}}
\newcommand\codv[3][]{\ThesisCompactDeriv{\mathrm{d}}{#1}{#2}{#3}}   % compact ordinary derivative: d_x y
\newcommand\cpdv[3][]{\ThesisCompactDeriv{\partial}{#1}{#2}{#3}}     % compact partial derivative: \partial_t u
\newcommand\cmdv[3][]{\ThesisCompactDeriv{\mathrm{D}}{#1}{#2}{#3}}   % compact material derivative: D_t u

\newcommand\transpose{^\textsf{T}}              % places in a transpose symbol

%% A command that LaTeX or the maths font already has (\div, \Re, \vec, ...) is redefined at the start
%% of the document (\AtBeginDocument): with MathTime, unicode-math sets these up only then, and would
%% undo a redefinition made earlier. Do the same for redefinitions of your own.
\newcommand{\grad}{\nabla}                     % gradient
\AtBeginDocument{\renewcommand{\div}{\nabla\cdot}} % divergence
\newcommand{\curl}{\nabla\times}               % curl

%% Dimensionless numbers of fluid mechanics and heat transfer (\mathword: upright, and read as one word,
%% "Re", by screen readers; accessibility.tex)
\AtBeginDocument{\renewcommand{\Pr}{\mathword{Pr}}} % Prandtl number
\AtBeginDocument{\renewcommand{\Re}{\mathword{Re}}} % Reynolds number
\newcommand{\Fr}{\mathword{Fr}}                   % Froude number
\newcommand{\Eu}{\mathword{Eu}}                   % Euler number
\newcommand{\St}{\mathword{St}}                   % Strouhal number
\newcommand{\Gr}{\mathword{Gr}}                   % Grashof number
\newcommand{\Ma}{\mathword{Ma}}                   % Mach number
\newcommand{\Pe}{\mathword{Pe}}                   % Péclet number
\newcommand{\Tu}{\mathword{Tu}}                   % Turbulence level

\DeclareMathOperator*{\argmax}{arg\,max}
\DeclareMathOperator*{\argmin}{arg\,min}

%% double dot product of tensors
\newcommand{\cddot}{\mathbin{{\cdot}{\cdot}}}

% L2 inner product, with the domain (which defines the inner product) as subscript:
%   \inner{\vec{v}}{\vec{u}}_\Omega            (v, u)_Omega, brackets at normal size (running text)
%   \inner[\big]{\vec{v}}{\vec{u}}_\Omega      larger brackets: \big, \Big, \bigg, \Bigg (displays)
%   \inner*{\vec{v}}{\frac{a}{b}}_\Omega       brackets sized to the contents
% Write time derivatives inside it as \cpdv{\vec{u}}{t} (gives \partial_t u): shorter than a fraction.
\DeclarePairedDelimiterX{\inner}[2]{(}{)}{#1, #2} % chktex 9 (Overleaf's checker misreads the ( and ) arguments)

%% Brackets that grow with their contents, all used like \inner: \norm{x}, \norm[\big]{x}, \norm*{\frac{a}{b}}
\DeclarePairedDelimiter{\norm}{\lVert}{\rVert}     % norm ||x||; with the space as subscript: \norm{\vec{u}}_{L^2}
\DeclarePairedDelimiter{\abs}{\lvert}{\rvert}      % absolute value |x|
\DeclarePairedDelimiter{\floor}{\lfloor}{\rfloor}  % floor
\DeclarePairedDelimiter{\ceil}{\lceil}{\rceil}     % ceiling
\DeclarePairedDelimiter{\jump}{\llbracket}{\rrbracket} % jump across a face [[u]] (discontinuous Galerkin)
\DeclarePairedDelimiterX{\avg}[1]{\{}{\}}{\mkern-4mu\delimsize\{#1\delimsize\}\mkern-4mu} % average {{u}}
\DeclarePairedDelimiterX{\dual}[2]{\langle}{\rangle}{#1, #2} % duality pairing <f, v>
\DeclarePairedDelimiter{\set}{\{}{\}}              % set by its elements: \set{1, 2, 3}
%% Set by a condition, the bar growing with the braces: \Set{x \in \R}{x > 0}, \Set*{\vec{v}}{\norm{\vec{v}} \leq 1}
\DeclarePairedDelimiterX{\Set}[2]{\{}{\}}{#1 \nonscript\:\delimsize\vert\allowbreak\nonscript\:\mathopen{} #2}
%% Restriction of a function to a set, f|_A: \restrict{f}{\Gamma}; \restrict[\big]{f}{\Gamma}; \restrict*{...}{\Gamma}
%% with the bar as tall as the function
\NewDocumentCommand{\restrict}{s o m m}{%
    \IfBooleanTF{#1}{\left.#3\right\rvert_{#4}}{\IfNoValueTF{#2}{#3\rvert_{#4}}{#3#2\rvert_{#4}}}}

%% Number sets, the constants e and i, upright as ISO 80000-2 sets them (\i outside maths stays the dotless i)
\newcommand{\R}{\mathbb{R}}                      % real numbers
\newcommand{\C}{\mathbb{C}}                      % complex numbers
\newcommand{\N}{\mathbb{N}}                      % natural numbers
\newcommand{\Z}{\mathbb{Z}}                      % integers
\newcommand{\e}{\mathrm{e}}                      % Euler's number: \e^{\i\omega t}
\makeatletter
\let\thesis@dotlessi\i
\DeclareRobustCommand{\i}{\ifmmode\mathrm{i}\else\thesis@dotlessi\fi} % imaginary unit
\makeatother

%% Operators, upright and spaced like \sin (\Re and \Im are the Reynolds number and free, see above)
\DeclareMathOperator{\tr}{tr}                     % trace
\DeclareMathOperator{\diag}{diag}                 % diagonal matrix
\DeclareMathOperator{\rank}{rank}                 % rank
\DeclareMathOperator{\sgn}{sgn}                   % sign
\DeclareMathOperator{\Real}{Re}                   % real part
\DeclareMathOperator{\Imag}{Im}                   % imaginary part
\DeclareMathOperator{\E}{\mathbb{E}}              % expectation: \E[X] or \E(X)
\DeclareMathOperator{\Var}{Var}                   % variance
\DeclareMathOperator{\Cov}{Cov}                   % covariance

%% Vectors and tensors in bold italics (Latin and Greek alike), so they stand apart from scalars,
%% e.g. the test function v from the viscosity nu. Change these two lines to switch notation everywhere:
%% plain italics: {#1}, arrow: {\vv{#1}} (package esvect), underline: {\underline{#1}}.
%% (\thesisboldmath is the bold italic of the maths font, set in fontMath.tex.)
\AtBeginDocument{\renewcommand{\vec}[1]{\thesisboldmath{#1}}} % vector:  \vec{u}
\newcommand{\tensor}[1]{\thesisboldmath{#1}}    % tensor:  \tensor{\tau}
